% Generated draft: requires mathematical and source review.
\documentclass[11pt]{article}
\usepackage[margin=25mm]{geometry}
\usepackage{fontspec}
\setmainfont{Latin Modern Roman}
\usepackage{amsmath,amssymb,mathtools,unicode-math}
\setmathfont{Latin Modern Math}
\usepackage{enumitem,needspace,xurl,hyperref}
\setlength{\parindent}{0pt}
\setlength{\parskip}{5pt}
\setlength{\emergencystretch}{3em}
\newcommand{\N}{\mathbb{N}}
\newcommand{\Z}{\mathbb{Z}}
\newcommand{\Q}{\mathbb{Q}}
\newcommand{\R}{\mathbb{R}}
\newcommand{\C}{\mathbb{C}}
\newcommand{\F}{\mathbb{F}}
\newcommand{\A}{\mathbb{A}}
\newcommand{\PP}{\mathbb{P}}
\newcommand{\E}{\mathbb{E}}
\newcommand{\eps}{\varepsilon}
\newcommand{\dd}{\,\mathrm{d}}
\begin{document}
\section*{UnsolvedMath ID: 800665}
\textbf{Source problem code:} catalogue entry
\subsection{Definitions and mathematical statement}
For a topological space \(X\), write \(\overline{U}^{\,X}\) for the closure of \(U\) in \(X\). Define
\[
\mathcal K(X)=\left\{\mathcal U\subseteq\{U\subseteq X:U\text{ is open in }X\}:\ X=\bigcup_{U\in\mathcal U}\overline{U}^{\,X}\right\}.
\]
Define \(\mathcal K_\Omega(X)\) to consist of those \(\mathcal U\in\mathcal K(X)\) such that
\[
(\forall U\in\mathcal U)\ \overline{U}^{\,X}\ne X,
\qquad
(\forall F\subseteq X\text{ finite})(\exists U\in\mathcal U)\ F\subseteq\overline{U}^{\,X}.
\]
Thus no member of \(\mathcal U\) is dense in \(X\). No countability restriction on these families is stated in the supplied record.

For a family \(\mathcal U\), let \([\mathcal U]^2\) denote its set of two-element subsets. The partition relation
\[
\mathcal K_\Omega(X)\longrightarrow\bigl(\mathcal K(X)\bigr)^2_2
\]
means that for every \(\mathcal U\in\mathcal K_\Omega(X)\) and every coloring \(c:[\mathcal U]^2\to\{0,1\}\), there exist \(\mathcal V\subseteq\mathcal U\) and \(i\in\{0,1\}\) such that \(\mathcal V\in\mathcal K(X)\) and
\[
(\forall p\in[\mathcal V]^2)\ c(p)=i.
\]

The supplied source reports that this partition relation implies that \(X\) is a Luzin set, attributing that implication to \([\mathrm{SchLuzin}]\). Its question, attributed to \([\mathrm{SchLuzin},\text{ Problem }4]\), is whether the converse holds, and hence whether
\[
X\text{ is a Luzin set}
\quad\Longleftrightarrow\quad
\mathcal K_\Omega(X)\longrightarrow\bigl(\mathcal K(X)\bigr)^2_2.
\]
This proposed equivalence must be read within the source's intended domain of spaces; that domain and the definition of ``Luzin set'' are not supplied. Under the usual real-line convention, a Luzin set is an uncountable set \(X\subseteq\R\) such that \(X\cap M\) is at most countable for every meager subset \(M\) of \(\R\), where meager means a countable union of nowhere dense sets. Under that convention, \(X\) would carry the subspace topology from \(\R\). This convention is recorded as a possible interpretation, not as an established hypothesis of the supplied source. The missing domain and convention must be resolved before the displayed equivalence is a fully specified assertion.
\subsection{Short English statement}
The question asks whether every Luzin set has the property that every two-coloring of pairs from a family in \(\mathcal K_\Omega(X)\) admits a monochromatic subfamily whose closures cover \(X\). The source reports the reverse implication. Its intended class of spaces and definition of Luzin set require clarification.
\subsection{Sources}
\begin{itemize}
\item Local UnsolvedMath registry, permanent ID \(800665\), source problem code \texttt{catalogue entry}; supplied title: ``Selection Principles Milestone Problems --- Problem 38 ([SchLuzin, Problem 4]).'' The registry attributes the record to Tsaban, with proposed year 2003, and labels it \emph{open}; this supplied status is not an independently verified current assessment.
\item Tsaban, \emph{Selection Principles in Mathematics: A Milestone of Open Problems} (2003), Problem 38. Supplied source URL: \url{https://arxiv.org/abs/math/0312182}. The record gives access date 2026-10-16 and extraction method \texttt{arxiv\_latex}; these are supplied provenance details, not a claim of independent access or verification.
\item The source attributes the reported implication to \([\mathrm{SchLuzin}]\) and the question to \([\mathrm{SchLuzin},\text{ Problem }4]\). No full bibliographic entry or separate URL for this reference is supplied.
\end{itemize}
\end{document}
